%%%%%%%%%%  scan idx 613  |  folio 600  %%%%%%%%%%
\SGAPrintPage{600}%
\[
\begin{tikzcd}
Y & \coprod_j Y_j \arrow[l]\\
X \arrow[u] & \coprod_j \Sigma_j \arrow[l] \arrow[u] & ,
\end{tikzcd}
\]
\[
\begin{tikzcd}
\underline{\Pic}_{Y/S} \arrow[r,"\sgacircled{1}"] \arrow[d,"\sgacircled{2}"'] &
\underline{\Pic}_{\coprod_jY_j/S}=\prod_j\underline{\Pic}_{Y_j/S} \arrow[d,"\sgacircled{4}"]\\
\underline{\Pic}_{X/S} \arrow[r,"\sgacircled{3}"] &
\underline{\Pic}_{\coprod_j\Sigma_j/S}=\prod_j\underline{\Pic}_{\Sigma_j/S}.
\end{tikzcd}
\]
And to establish II, it suffices to see that, after shrinking \(S\), \(\sgacircled{1}\), \(\sgacircled{3}\) and \(\sgacircled{4}\) are quasi-affine.

Consider \(\sgacircled{4}\) first. Replacing \(Y\) by \(Y_j\) and \(X\) by \(\Sigma_j\) reduces us to the case where \(Y\) is integral over \(S\). Since \(f\) is surjective, one of the \(X_i\), say \(X_0\), dominates \(Y\), and we have the two sequences
\[
  X_0\hookrightarrow \coprod X_i\longrightarrow Y,
\]
\[
  \Pic_{X_0/S}\longleftarrow \Pic_{\coprod X_i/S}\longleftarrow \Pic_{Y/S}.
\]
Thus, knowing that the \(\underline{\Pic}_{X_i/S}\), and therefore also \(\Pic_{\coprod X_i/S}\), are representable and that the composite morphism is quasi-affine, we conclude that \(\sgacircled{4}\) is quasi-affine. In view of the proof of 1.2, for \(\sgacircled{4}\) we are therefore reduced to proving II'. \uline{Theorem} 1.1 \uline{when} \(X\) \uline{and} \(Y\) \uline{are integral over} \(S\).

For \(\sgacircled{1}\) and \(\sgacircled{3}\), we must treat the case of a morphism \(\coprod X_i\to X\) where \(X\) is reduced, the \(X_i\) are the irreducible components of \(X\) and are integral over \(S\). Now we may suppose that \(X\) has connected fibers. Indeed, let \(\eta\) be the geometric point of \(S\). Then the irreducible components of \(X_\eta\) are the \((X_i)_\eta\) and, after shrinking \(S\), we may suppose that \(X_i\cap X_j=\varnothing\) if and only if


%%%%%%%%%%  scan idx 614  |  folio 601  %%%%%%%%%%
\SGAPrintPage{601}%
\((X_i)_\eta\cap(X_j)_\eta=\varnothing\). Let \(Y_k=\bigcup_{i\in I(k)}X_i\) be the connected components of \(X\); then \(f\) decomposes, \(f=\coprod_k f_k:\coprod_k(\coprod_{i\in I(k)}X_i)\longrightarrow \coprod_kY_k\), and the assertion follows. Finally, an immediate induction reduces us to the case of a morphism \(X_1\amalg X_2\to X\). After shrinking \(S\), \(\sgacircled{1}\) and \(\sgacircled{3}\) then follow, by (EGA III 7.8.6), from the

\sgaStableTarget{sga6.e12.prop.2-3}\textbf{Proposition 2.3.} \uline{Let} \(S\) \uline{be a scheme,} \(g:Y\to S\) \uline{a proper morphism, and} \(Y_1,Y_2\) \uline{two closed subschemes defined by the ideals} \(I_1,I_2\)\uline{. Suppose that}
\begin{enumerate}[label=\alph*)]
\item \(Y,Y_1,Y_2,Y_1\cap Y_2\) \uline{are faithfully flat and of finite presentation over} \(S\).
\item \(g_*{O}_Y={O}_S\) and \((g_i)_*{O}_{Y_i}={O}_S\) \uline{universally (where} \(g_i=g|Y_i\), \(i=1,2\)).
\item \(I_1\cap I_2=0\).
\end{enumerate}
\uline{Then} \(\underline{\Pic}_{Y/S}\longrightarrow \underline{\Pic}_{(Y_1\amalg Y_2)/S}\) \uline{is representable by an affine morphism of finite presentation.}

Indeed, first note that the hypotheses are stable under every base change. For (c), this follows from the two exact sequences:
\[
  0\longrightarrow I_1+I_2\longrightarrow{O}_Y
  \longrightarrow{O}_{Y_1\cap Y_2}\longrightarrow0
\]
and
\[
  0\longrightarrow I_1\cap I_2\longrightarrow I_1\oplus I_2
  \longrightarrow I_1+I_2\longrightarrow0~;
\]
the first shows that \(I_1+I_2\) is flat over \(S\); then the second shows that \(I_1\cap I_2\) is stable under base change.

The assertion is local for the f.p.p.f. topology. Making the base change \(Y_1\cap Y_2\to S\), we may therefore suppose that \(Y_1\cap Y_2/S\) has a section \(\alpha\), which defines sections \(\alpha_i\) of \(Y_i\) and \(\beta\) of \(Y\).

After a suitable base change, we are reduced to seeing that every point \(\mu:S\to\underline{\Pic}_{(Y_1\amalg Y_2)/S}\) has as its inverse image a subfunctor of


%%%%%%%%%%  scan idx 615  |  folio 602  %%%%%%%%%%
\SGAPrintPage{602}%
\(\underline{\Pic}_{Y/S}\) that is representable by an affine scheme over \(S\). Now \(\mu\) is defined by an invertible sheaf \(L_i\) on \(Y_i\) equipped with a rigidification \(u_i:\alpha_i^*L_i\xrightarrow{\sim}{O}_S\) for \(i=1,2\). Likewise, a point \(\lambda:S\to\underline{\Pic}_{Y/S}\) corresponds to an invertible sheaf \(L\) on \(Y\) equipped with a rigidification \(u:\alpha^*L\xrightarrow{\sim}{O}_S\). The image of \(\lambda\) in \(\Pic_{Y_1\amalg Y_2/S}\) is \(\mu\) if and only if there is an isomorphism of rigidified sheaves \(f_i^*(L)\simeq L_i\) where \(f_i:Y_i\hookrightarrow Y\) is the inclusion for \(i=1,2\). Observe that these isomorphisms define an isomorphism \(\varphi:L_1|Y_1\cap Y_2\xrightarrow{\sim}L_2|Y_1\cap Y_2\) such that \(u_1=u_2\circ\varphi\).

Conversely, given an isomorphism \(\varphi\), \(L_1|Y_1\cap Y_2\to L_2|Y_1\cap Y_2\), such that \(u_1=u_2\circ\varphi\), there is a rigidified invertible sheaf \(L\) on \(Y\) obtained by gluing along \(\varphi\); explicitly, \(L\) is the kernel of the canonical homomorphism \((f_1)_*L_1\oplus(f_2)_*L_2\longrightarrow(f|Y_1\cap Y_2)_*L_2\) defined by \(\varphi\). This \(L\) then corresponds to a point \(\lambda\) of \(\underline{\Pic}_{Y/S}\) mapping to \(\mu\).

Finally, let \(I\) be the functor defined, as \(T\) runs through the schemes over \(S\), by the sets of isomorphisms \(\varphi:(L_1)_T|(Y_1\cap Y_2)_T\xrightarrow{\sim}(L_2)_T|(Y_1\cap Y_2)_T\) such that \((u_1)_T=(u_2)_T\circ\varphi\). We have just established an isomorphism between the inverse image of \(\mu\) in \(\underline{\Pic}_{Y/S}\) and the functor \(I\). It therefore suffices to apply the following lemma (taking \(X=Y_1\cap Y_2\)):

\sgaStableTarget{sga6.e12.lem.2-4}\textbf{Lemma 2.4.} \uline{Let} \(X\) \uline{be a flat, proper scheme of finite presentation over} \(S\)\uline{. Let} \(L,M\) \uline{be two invertible sheaves on} \(X\), \uline{and} \(\underline{\Isom}(L,M)\) \uline{the functor of isomorphisms} \(L_T\to M_T\) \uline{(for} \(T/S\)\uline{). Then}
\begin{enumerate}[label=\roman*)]
\item \(\underline{\Isom}(L,M)\) \uline{is representable by an affine scheme over} \(S\) \uline{of finite presentation.}
\item \uline{If, moreover,} \(L\) \uline{and} \(M\) \uline{are equipped with rigidifications (relative to a}


%%%%%%%%%%  scan idx 616  |  folio 603  %%%%%%%%%%
\SGAPrintPage{603}%
\uline{given section of} \(X\)\uline{), the subfunctor of} \(\underline{\Isom}(L,M)\) \uline{formed by the isomorphisms compatible with these rigidifications is representable by a closed subscheme of} \(\underline{\Isom}(L,M)\).
\end{enumerate}

Indeed, ii) is clear. As for i), recall that the functors \(T\mapsto\Hom(L_T,M_T)\) and \(T\mapsto\Hom(M_T,L_T)\) are representable by vector bundles \(\mathbb V(E)\) and \(\mathbb V(F)\) (by EGA III 7.7.8). Then, for \(u\in\Hom(L_T,M_T)\) and \(v\in\Hom(M_T,L_T)\), the conditions \(u\circ v=1\) and \(v\circ u=1\) define a closed subscheme \(F\) of \(\mathbb V(E)\times_S\mathbb V(F)\), and it is clear that \(F\simeq I\).

\sgaStableTarget{sga6.e12.para.2-5}\textbf{2.5.} Finally, let us reduce assertion II' to the following assertion:

\(\mathrm{II}^*\). \uline{Theorem} 1.1 \uline{when} \(X\) \uline{and} \(Y\) \uline{are integral over} \(S\) \uline{and the sequence}
\sgaStableTarget{sga6.e12.eq.unnamed}\begin{equation}\tag{*}
  {O}_Y \longrightarrow f_*{O}_X \rightrightarrows g_*\bigl({O}_{X\times_YX}\bigr)
\end{equation}
\uline{(where} \(g:X\times_YX\to Y\) \uline{is the structural morphism) is exact.}

In fact, for this it suffices to prove that there is a factorization of \(f\)
\[
  X\xrightarrow{f_0}Y_0\xrightarrow{f_1}Y_1\longrightarrow\cdots\xrightarrow{f_\eta}Y_\eta=Y
\]
such that every \(f_i\) satisfies (*) and, after shrinking \(S\), every \(Y_i\) is integral over \(S\).

For this, put \(Y_0=\Spec(f_*(\underline{O}_X))\), with \(f_0\) arising from the Stein factorization of \(f\) (EGA III 4.3.3), so that \(f_{0*}(\underline{O}_X)\xleftarrow{\sim}\underline{O}_{Y_0}\), and hence \(f_0\) does satisfy (*). To factor \(Y_0\to Y\) in the same way, we use the

\sgaStableTarget{sga6.e12.lem.2-6}\textbf{Lemma 2.6.} \uline{Let} \(Y\) \uline{be a noetherian scheme, and} \(f:X\to Y\) \uline{a finite epimorphism (i.e. such that} \({O}_Y\to f_*({O}_X)\) \uline{is injective); then} \(f\) \uline{factors into a finite sequence of finite morphisms} \(f_i\) \(X\simeq Y_0\xrightarrow{f_1}Y_1\longrightarrow\cdots\longrightarrow Y_n=Y\) \uline{that are effective epimorphisms (i.e. such that the sequence}


%%%%%%%%%%  scan idx 617  |  folio 604  %%%%%%%%%%
\SGAPrintPage{604}%
\({O}_{Y_i}\longrightarrow (f_i)_*({O}_{Y_{i-1}})\rightrightarrows (g_i)_*\bigl({O}_{Y_{i-1}}\otimes_{{O}_{Y_i}}{O}_{Y_{i-1}}\bigr)\) \uline{is exact).}

Put \(A={O}_Y\), \(B_0=f_*({O}_X)\), and define \(B_i=\Ker\bigl(B_{i-1}\rightrightarrows B_{i-1}\otimes_A B_{i-1}\bigr)\) recursively. We verify at once that \(B_{i-1}\otimes_A B_{i-1}\xrightarrow{\sim}B_{i-1}\otimes_{B_i}B_{i-1}\). It remains to prove that, for \(i\) sufficiently large, \(B_i=A\). Indeed, consider the decreasing sequence of \(A\)-modules \(M_i=B_i/A\), and put \(S_i=\Supp(M_i)\). The decreasing sequence of closed subsets \(S_i\) of \(Y\) is stationary; let \(S\) be the stationary value. Let us show that \(S=\varnothing\). Otherwise, let \(s\) be a maximal point of \(S\). After replacing \(Y\) by \(\Spec{O}_{Y,s}\), we may suppose that \(Y\) is local with closed point \(s\). By the choice of \(s\), \(M_i\) is artinian for \(i\) sufficiently large, so the sequence \((M_i)\) is stationary. Thus there exists \(i\) such that \(B_i=B_{i+1}\). Consequently, \(\Spec(B_i)\to Y\) is a monomorphism, hence a closed immersion (Nakayama); since \(A\subset B_i\), we have \(A=B_i\) and \(s\notin S\), a contradiction.

\sgaStableTarget{sga6.e12.sec.3}\section*{3. Proof of I: Oort's dévissage}

\sgaStableTarget{sga6.e12.prop.3-1}\textbf{Proposition 3.1.} (Oort, Bull. soc. math. France, vol. 90-1962). \uline{Let} \(S\) \uline{be a locally noetherian scheme,} \(f:X\to S\) \uline{a proper morphism, and} \(X'\subset X\) \uline{a closed subscheme. Suppose that:}
\begin{enumerate}[label=\roman*)]
\item \(X'\) \uline{is defined by an ideal} \(I\) \uline{contained in the nilradical} \(N\) \uline{of} \(X\) \uline{and} \(N.I=0\).
\item \(X'\) \uline{is flat over} \(S\) \uline{and cohomologically flat over} \(S\) \uline{in dimension} \(0\) (EGA III 7.8.1).
\item \(I\) \uline{is flat over} \(S\) \uline{and cohomologically flat over} \(S\) \uline{in dimensions} \(0\) \uline{and} \(1\)\uline{, so that (by} EGA III 7.4.6\uline{) there are a coherent} \({O}_S\)\uline{-module} \(Q\) \uline{and an isomorphism, functorial in the coherent} \({O}_S\)\uline{-module} \(M\),

